\documentclass[a4paper,zihao=-4,fontset=fandol]{ctexart}
\usepackage[left=25mm,right=25mm,top=25mm,bottom=25mm,footskip=12mm]{geometry}
\usepackage{amsmath,amssymb,bm,graphicx,booktabs,tabularx,longtable,array,float}
\usepackage{caption,fancyvrb,fvextra,algorithm,algorithmic,hyperref}
\usepackage{xurl,tikz,pdfpages,placeins}
\usetikzlibrary{arrows.meta,positioning,calc}
\hypersetup{colorlinks=true,linkcolor=blue,citecolor=blue,urlcolor=blue,filecolor=blue,menucolor=blue,runcolor=blue,pdftitle={张量驻留与查询时序驱动的神经网络处理器多核调度},pdfauthor={},pdfsubject={2026 A题论文初稿}}
\setmainfont{texgyretermes}[Extension=.otf,UprightFont=*-regular,BoldFont=*-bold,ItalicFont=*-italic,BoldItalicFont=*-bolditalic]
% Derived from the D-template font configuration; portable font detection.
\ifcontest
  \IfFileExists{/Applications/Microsoft Word.app/Contents/Resources/DFonts/Simsun.ttc}{
    \setCJKmainfont[AutoFakeBold=2,Path={/Applications/Microsoft Word.app/Contents/Resources/DFonts/}]{Simsun.ttc}
    \setCJKsansfont[AutoFakeBold=2,Path={/Applications/Microsoft Word.app/Contents/Resources/DFonts/}]{SimHei.ttf}
    \setCJKmonofont[Path={/Applications/Microsoft Word.app/Contents/Resources/DFonts/}]{Simsun.ttc}
  }{\IfFontExistsTF{SimSun}{
    \setCJKmainfont[AutoFakeBold=2]{SimSun}\setCJKsansfont[AutoFakeBold=2]{SimHei}\setCJKmonofont{SimSun}
  }{\setCJKmainfont{FandolSong-Regular}\setCJKsansfont{FandolHei-Regular}\setCJKmonofont{FandolSong-Regular}
    \typeout{A-PAPER FONT MODE: Fandol preview; official Song/Hei fonts are unavailable.}
  }}
\fi
\newcommand{\xinweifont}{\kaishu}

\ifcontest
  \linespread{1.0}
\else
  \geometry{top=30mm}
  \linespread{1.38}
\fi
\setlength{\parindent}{2em}
\setlength{\parskip}{0pt}
\setlength{\emergencystretch}{2em}
\ctexset{section={name={,、},number=\chinese{section},aftername={},
  format=\centering\sffamily\zihao{4},beforeskip=18bp,afterskip=12.4bp},
subsection={format=\zihao{-4},titleformat=\sffamily,beforeskip=12bp,afterskip=6.8bp},
subsubsection={format=\zihao{-4},titleformat=\sffamily,beforeskip=10bp,afterskip=6bp}}
\pagestyle{plain}
\renewcommand{\thefigure}{\arabic{figure}}
\renewcommand{\thetable}{\arabic{table}}
\renewcommand{\theequation}{\arabic{equation}}
\captionsetup{font=normalsize,labelsep=quad,skip=8bp}
\captionsetup[figure]{font={small,bf},labelsep=space,position=bottom,skip=4bp}
\captionsetup[table]{font={small,bf},labelsep=space,position=top,skip=4bp}
\floatname{algorithm}{算法}
\renewcommand{\algorithmicrequire}{\textbf{输入：}}
\renewcommand{\algorithmicensure}{\textbf{输出：}}
\newcolumntype{Y}{>{\raggedright\arraybackslash}X}
\newcommand{\pending}{\textbf{待计算}}
\newcommand{\slot}[1]{\textbf{【#1】}}
\newcommand{\E}{\mathbb{E}}
\newcommand{\R}{\mathbb{R}}
\newcommand{\samplefigure}[3]{%
  \begin{figure}[H]\centering
  \includegraphics[width=0.91\linewidth,height=0.43\textheight,keepaspectratio]{figures/#1.pdf}
  \caption{#2}\label{#3}\end{figure}}
\newcounter{extrafigure}
\renewcommand{\theextrafigure}{S.\arabic{extrafigure}}
\newcommand{\auxfigure}[3]{%
  \begin{figure}[H]\centering
  \includegraphics[width=0.91\linewidth,height=0.40\textheight,keepaspectratio]{figures/#1.pdf}
  \refstepcounter{extrafigure}\caption*{图\theextrafigure\quad #2}\label{#3}\end{figure}}
\newcommand{\auxtablecaption}[1]{\caption{#1}}
\newcommand{\aitool}{OpenAI Codex}
\newcommand{\aimodel}{GPT-6（本文修订）；原实验精确型号待核实}
\newcommand{\aideveloper}{OpenAI}
\newcommand{\aireleased}{待核实}
\newcommand{\Case}[1]{\texttt{case\_#1}}
\tikzset{abox/.style={draw,align=center,text width=3.5cm,minimum height=1cm,inner sep=5pt},alink/.style={-{Stealth[length=2mm]},thick}}
\newcommand{\modelcard}[4]{\subsection{#1}#2\par\textbf{基本形式：}#3\par\textbf{选择与检验：}#4\par}
